\documentclass[11pt,a4paper]{article}

\usepackage[utf8]{inputenc}
\usepackage[T1]{fontenc}
\usepackage{amsmath, amssymb}
\usepackage{graphicx}
\usepackage{hyperref}
\usepackage[authoryear]{natbib}
\usepackage{geometry}
\usepackage{microtype}
\usepackage{xcolor}
\usepackage{booktabs}
\usepackage{enumitem}
\usepackage{tcolorbox}

\geometry{margin=2.5cm}

\hypersetup{
  colorlinks=true,
  linkcolor=blue!60!black,
  citecolor=blue!60!black,
  urlcolor=blue!60!black,
}

% AAS journal-abbreviation macros (not loading aastex, so define them).
\providecommand{\mnras}{MNRAS}
\providecommand{\mnrasl}{MNRAS Letters}
\providecommand{\apj}{ApJ}
\providecommand{\apjs}{ApJS}
\providecommand{\apjl}{ApJL}
\providecommand{\aj}{AJ}
\providecommand{\aap}{A\&A}
\providecommand{\jcap}{JCAP}
\providecommand{\prd}{Phys.~Rev.~D}

\newcommand{\sigmaoff}{\sigma_{\rm off}}
\newcommand{\fcen}{f_{\rm cen}}
\newcommand{\dsig}{\Delta\Sigma}
\newcommand{\dm}{\Delta m}
\newcommand{\Rmin}{R_{\rm min,fit}}

\title{\textbf{Combined Recovery of the Mass--Concentration Relation and
Magnitude-Gap Slope from Un-Stacked DES~Y3-like Cluster Lensing Profiles}\\[0.4em]
\large A description of \texttt{combined\_recovery\_lensing\_triaxial\_unbiased.ipynb}}
\author{Andr\'e Vitorelli}
\date{May 2026}

\begin{document}

\maketitle

\begin{abstract}
We describe the end-to-end mock pipeline implemented in
\texttt{notebooks/combined\_recovery\_lensing\_triaxial\_unbiased.ipynb}, the
current most-advanced experiment in the \textsc{destacking} package.  The
pipeline (i)~draws $N=16\,000$ mock clusters with masses from the Tinker08
halo mass function and DES~Y3-like \citep{Costanzi2021} redshifts, (ii)~assigns
triaxial NFW shapes with CDM-like axis ratios and isotropic line-of-sight
orientations, (iii)~simulates the lensing signal of a metacalibration-format
source field with a DES~Y3-like Smail--Schechter $n(z)$, (iv)~applies a standard
Rayleigh + delta-function miscentering model with $\fcen=0.64$,
$\sigmaoff=0.32\,h^{-1}\mathrm{Mpc}$, and (v)~jointly fits the
$(c_0, \alpha, \beta, \gamma, \log M_{\rm norm})$ scaling-relation parameters
under a spherical-NFW forward model with fixed-parameter miscentering and a
radial cut at $R \ge 0.5\,\mathrm{Mpc}$.  Priors are tight Gaussians centred on
the Dutton \& Maccio (2014) power-law expectation, except for $\gamma$ which is
$\mathcal{N}(0, 1)$.  We discuss the diagnostic that motivates the radial cut,
the role of each fixed parameter, and a list of concrete steps required before
the pipeline is applied to real DES~Y3 data.
\end{abstract}

\tableofcontents
\newpage

%──────────────────────────────────────────────────────────────────────────────
\section{Goals}
%──────────────────────────────────────────────────────────────────────────────

The mass--concentration ($M$--$c$) relation of dark-matter halos has been
shown in $N$-body simulations to be a function of mass, redshift, and a few
secondary halo properties associated with formation history
\citep{Wechsler2002, DuttonMaccio2014, Diemer2019}.  Among optical
observables, the \emph{magnitude gap}
$\dm = m_2 - m_1$ between the brightest cluster galaxy (BCG) and the
brightest satellite within $0.5\,R_{200}$ is a proxy for assembly age
\citep{Jones2003, Khosroshahi2006}: clusters that have grown few satellites
since the BCG dominated their light tend to be older systems.

This notebook tests whether the destacking technique—joint Gaussian likelihood
fit of $N$ individual per-cluster $\dsig$ profiles
\citep{Vitorelli2025}—can simultaneously recover the baseline
$M$--$c$ relation and the magnitude-gap slope
\begin{equation}
  \gamma \equiv \frac{\mathrm{d}\log_{10}c}{\mathrm{d}\log_{10}(\dm/\dm_{\rm ref})},
  \label{eq:gamma_def}
\end{equation}
from data simulated to mimic the DES~Y3 cluster lensing environment, including:
triaxial halo geometry, a realistic source $n(z)$, and a Rayleigh-distributed
miscentering offset.

%──────────────────────────────────────────────────────────────────────────────
\section{Simulation pipeline}
%──────────────────────────────────────────────────────────────────────────────

\subsection{Parameters}

The headline configuration in the notebook is summarised in
Table~\ref{tab:params}.

\begin{table}[h]
\centering
\begin{tabular}{lll}
\toprule
Parameter & Symbol & Value \\
\midrule
Number of clusters             & $N$               & 16\,000 \\
True magnitude-gap slope       & $\gamma_{\rm true}$ & 0.2 \\
Reference magnitude gap        & $\dm_{\rm ref}$     & 0.1\,mag \\
Radial bins                    & RADII             & $10^{-0.8\text{--}+0.8}$\,Mpc, 12 log-spaced \\
Radial cut                     & $\Rmin$           & 0.5\,Mpc \\
\midrule
Cluster $n(z)$ range           & $[z_{\rm min}, z_{\rm max}]$ & $[0.2, 0.65]$ \\
Cluster $n(z)$ shape           &                   & $P(z) \propto z^2$ \\
Source $n(z)$ shape            &                   & Smail--Schechter \\
Source $n(z)$ parameters       & $(\alpha, \beta, z_0)$ & $(2.0, 1.5, 0.4)$ \\
Source density                 &                   & $\sim\!30\,\mathrm{gal\,arcmin}^{-2}$ \\
Source field half-width        &                   & 0.006\,rad ($\sim\!20'$) \\
Shape-noise per component      & $\sigma_e$        & 0.27 \\
\midrule
Centred fraction               & $\fcen$           & 0.64 \\
Miscentering Rayleigh scale    & $\sigmaoff$       & $0.32\,h^{-1}\mathrm{Mpc}$ \\
\midrule
Nautilus live points           & $n_{\rm live}$    & 400 \\
Parallel workers               & --                & 15 \\
\bottomrule
\end{tabular}
\caption{Top-level configuration of the unbiased combined-recovery notebook.}
\label{tab:params}
\end{table}

\subsection{Truth catalogue with a DES~Y3-like cluster $n(z)$}

The cluster truth catalogue is constructed in two stages:

\paragraph{Uniform-$z$ pool.} \texttt{generate\_truth\_catalog} draws an
oversampled pool of $5N$ clusters with redshifts uniform in $[z_{\rm min},
z_{\rm max}]$ and, at each redshift, halo masses sampled from the Tinker08
halo mass function \citep{Tinker2008} over
$\log_{10}M \in [13.5, 15]$.  Each cluster also receives:
\begin{itemize}[noitemsep]
  \item A BCG mass from
        $\log_{10}M_{\rm BCG} = 0.4\log_{10}M + 6 + \mathcal{N}(0, 0.15)$
        and luminosity $L_{\rm BCG}$ from the \citet{Zheng2007} BCG-mass
        relation.
  \item A satellite count drawn from the \citet{Zheng2007} HOD.
  \item Satellite luminosities from $\phi(L) \propto L^\alpha$ with
        $\alpha = -1.2$ on $[L_{\rm min}, L_{\rm BCG}]$.
  \item Magnitude gap $\dm = 2.5\log_{10}(L_{\rm BCG}/L_{\rm sat,max})$.
\end{itemize}
The baseline concentration is the \citet{DuttonMaccio2014} (DM14) prediction
\begin{align}
  \log_{10} c_{\rm DM14} &= a(z) + b(z)\log_{10}\!\left(\frac{M_{200m}}{10^{12}\,M_\odot}\right),
  \label{eq:DM14}\\
  a(z) &= 0.52 + 0.385\exp(-0.617\, z^{1.21}),\nonumber\\
  b(z) &= -0.101 + 0.026\,z, \nonumber
\end{align}
and the truth catalogue then imprints the magnitude-gap dependence as a
\emph{direct} $\log_{10}c$ shift:
\begin{equation}
  \log_{10} c_{\rm truth}
  = \log_{10} c_{\rm DM14}(z, M) + \gamma_{\rm true}\,\log_{10}\!\frac{\dm}{\dm_{\rm ref}}.
  \label{eq:c_truth}
\end{equation}

\paragraph{DES~Y3 reweighting.} The pool is then resampled without replacement
using the DES~Y3-like cluster $n(z)$:
\begin{equation}
  P_{\rm DESY3}(z) \propto z^2,\quad z\in[0.2, 0.65],
  \label{eq:cluster_nz}
\end{equation}
which is a volume-weighted approximation of the redMaPPer Y3 catalogue
\citep{Rykoff2016, Costanzi2021}.  Because resampling preserves each pool
cluster's $(z, M, c, \dm)$ entry, the mass--redshift correlation imposed by the
Tinker08 HMF is retained.  The final OLS fit
\begin{equation}
  \log_{10}c = \log_{10}c_0 + \alpha(\log_{10}M - 14) + \beta\log_{10}\!\frac{1+z}{1+z_{\rm ref}}
             + \gamma\log_{10}\!\frac{\dm}{\dm_{\rm ref}}
  \label{eq:parametric_c}
\end{equation}
of the truth concentrations gives the sample-specific
$(c_0, \alpha, \beta, \gamma)$ values used as the ``truth'' against which the
posterior is compared.

\paragraph{Source $n(z)$.}  The source $n(z)$ is the Smail--Schechter form
\citep{Smail1994}
\begin{equation}
  n(z) \propto (z/z_0)^\alpha \exp\!\left[-(z/z_0)^\beta\right],
  \label{eq:source_nz}
\end{equation}
with $\alpha=2$, $\beta=1.5$, $z_0=0.4$, peaking near $z\approx0.4$ and with
mean $\bar z \approx 0.66$ and a tail to $z\approx 2.5$—an approximation of
the DES~Y3 metacalibration source distribution
\citep{Myles2021, MacCrann2021}.  The sampler is called inside the lensing
worker to overwrite the uniform-$z$ default column produced by
\texttt{xlensing.sims.make\_simple\_random\_cat}.

\begin{figure}[ht]
  \centering
  \includegraphics[width=0.98\textwidth]{notebooks/des_y3_nz}
  \caption{Adopted DES~Y3-like redshift distributions used throughout the
  pipeline.  \emph{Left}: cluster $n(z) \propto z^2$ on $[0.2, 0.65]$, the
  volume-weighted approximation of the redMaPPer Y3 catalogue.  \emph{Right}:
  source $n(z)$ in the Smail--Schechter parametrisation with $(\alpha, \beta,
  z_0) = (2, 1.5, 0.4)$, an approximation of the DES~Y3 metacalibration source
  distribution.}
  \label{fig:nz}
\end{figure}

\paragraph{Catalogue distributions.}  After resampling, the marginal
distributions of cluster redshift, halo mass, concentration, and magnitude gap
are shown in Figure~\ref{fig:truth}.  Figure~\ref{fig:mc_scatter} shows the
mass--concentration plane coloured by $\dm$ alongside the magnitude-gap--vs--$c$
plane coloured by $\log_{10}M$; the visible tilt in the right panel is the
$\gamma$-induced signal that the inference is intended to recover.

\begin{figure}[ht]
  \centering
  \includegraphics[width=0.78\textwidth]{notebooks/truth_distributions}
  \caption{Marginal distributions of the simulated truth catalogue:
  cluster redshift (DES~Y3-like), halo mass (Tinker08 HMF restricted to
  $\log_{10}M\in[13.5, 15]$, log scale on the right), concentration after the
  $\gamma$ shift (Eq.~\ref{eq:c_truth}), and magnitude gap.}
  \label{fig:truth}
\end{figure}

\begin{figure}[ht]
  \centering
  \includegraphics[width=0.98\textwidth]{notebooks/mc_scatter}
  \caption{\emph{Left}: $\log_{10}M$ vs $c_{200}$ coloured by $\dm$, with the
  DM14 reference curve at the sample median redshift overlaid (black).
  \emph{Right}: magnitude gap vs $c_{200}$ coloured by $\log_{10}M$; the
  upward trend with $\dm$ is the $\gamma$-induced fossil-cluster signal.
  The title lists the OLS-extracted $(c_0, \alpha, \beta, \gamma)$ truth values
  for the sample.}
  \label{fig:mc_scatter}
\end{figure}

\subsection{Triaxial halo shape model}

Each cluster receives an axis-ratio realisation
$b/a \sim \mathcal{U}(0.5, 1)$, $c/a \sim \mathcal{U}(0.3, b/a)$, mimicking
CDM $N$-body shape distributions \citep{Allgood2006, Schneider2012}.
Line-of-sight directions are isotropic
($\cos\theta\sim\mathcal{U}(0,1)$, $\phi\sim\mathcal{U}(0,2\pi)$).  The
projected ellipse is computed via \texttt{xlensing.sims.project\_ellipsoid}
(Schur complement of the shape matrix), returning $(q_\mathrm{maj},
q_\mathrm{min}, \mathrm{PA})$.  The shear is then applied with
\texttt{xlensing.sims.apply\_triaxial\_NFW\_shear\_region} \citep{xlensing}.

\begin{figure}[ht]
  \centering
  \includegraphics[width=0.95\textwidth]{notebooks/triaxial_shapes}
  \caption{Triaxial-shape distributions.  Input 3-D axis ratios $b/a$ and
  $c/a$ (left, centre, blue) and the resulting projected axis ratio
  $q_{\rm min}/q_{\rm maj}$ on the sky (right, red).  The projected axis
  ratio peaks near $0.7$--$0.8$, consistent with the X-ray cluster morphology
  literature \citep{Kawahara2010, Sereno2018}.}
  \label{fig:triaxial}
\end{figure}

\subsection{Miscentering: simulation side}
\label{sec:miscentering_sim}

A fraction $\fcen=0.64$ of clusters are correctly centred and the remaining
$36\%$ are miscentred with a Rayleigh-distributed offset:
\begin{equation}
  P(R_{\rm mis}) =
    \fcen\,\delta(R_{\rm mis}) + (1-\fcen)\,
    \mathcal{R}(R_{\rm mis};\sigmaoff),
  \label{eq:mis_pdf}
\end{equation}
\begin{equation}
  \mathcal{R}(R;\sigma) = \frac{R}{\sigma^2}\exp\!\left(-\frac{R^2}{2\sigma^2}\right),\quad R\ge0.
  \label{eq:rayleigh}
\end{equation}
The default Rayleigh scale is $\sigmaoff = 0.32\,h^{-1}\,\mathrm{Mpc}$,
consistent with redMaPPer miscentering calibrations
\citep{Rykoff2016, Zhang2019}.  Each miscentred cluster receives a specific
offset magnitude and a uniformly random azimuth.  In the lensing worker the
triaxial NFW shear is applied around the \emph{true} centre $(0, 0)$, while
\texttt{xlensing.data.metacal\_cluster\_lensing} is called with the
\emph{displaced} centre
$(\Delta\mathrm{RA}, \Delta\mathrm{Dec}) = (R_{\rm mis}/D_A)(\cos\phi,\sin\phi)$.

\begin{figure}[ht]
  \centering
  \includegraphics[width=0.95\textwidth]{notebooks/miscentering_offsets}
  \caption{Realised miscentering offsets in the simulated catalogue (the
  miscentred subset only, i.e. the $36\%$ for which $R_{\rm mis} > 0$).
  \emph{Left}: histogram of the offset magnitude in physical Mpc, with the
  target Rayleigh PDF overlaid in solid black.  \emph{Right}: 2-D scatter of
  the offset positions on the sky relative to the true centre; the dashed
  circle has radius $\sigmaoff = 0.32\,h^{-1}\,\mathrm{Mpc} \approx 0.47$\,Mpc
  for the Planck18 $h$.}
  \label{fig:mis}
\end{figure}

\subsection{Lensing measurement}

A single background catalogue is built with
\texttt{make\_simple\_random\_cat}\,(density$\sim$30\,gal/arcmin$^2$,
half-width 0.006\,rad), its $z$ column is overwritten with samples from
Eq.~(\ref{eq:source_nz}), and metacalibration-format responses are assigned
from a lognormal size distribution.  Per cluster, intrinsic ellipticities are
drawn from \texttt{ngmix.GPriorBA} with $\sigma_e = 0.27$, the triaxial shear
is applied, sources are kept above $z_s > 1.1\,z_l + 0.2$, and
\texttt{xlensing.data.signal} bins the lensing signal into log-spaced annuli.
The covariance is diagonal in radial bins, with the bin-$b$ variance equal to
the per-galaxy shape-noise contribution
$\sigma_b^2 = (\sigma_e\,\bar\Sigma_{\rm cr,b}/\sqrt{n_b})^2$
when $n_b \ge 5$; empty or sparse bins are inflated to a noise floor
$\sigma_{\rm floor}^2 = 10^4\,(M_\odot/\mathrm{pc}^2)^2$ so they are effectively
excluded from the likelihood.

%──────────────────────────────────────────────────────────────────────────────
\section{Forward model and analysis}
%──────────────────────────────────────────────────────────────────────────────

\subsection{Joint Gaussian likelihood}

For each cluster $i$ let
$\mathbf{d}_i\in\mathbb{R}^{N_{\rm bin}}$ be the measured profile and
$\mathbf{m}_i(\boldsymbol\theta)$ the model prediction under parameters
$\boldsymbol\theta = (c_0, \alpha, \beta, \gamma, \log M_{\rm norm})$.  The
joint log-likelihood is
\begin{equation}
  \ln\mathcal{L}(\boldsymbol\theta) = -\frac{1}{2}\sum_{i=1}^N\!\left[
    k\ln(2\pi) + \ln\det\mathbf{C}_i + (\mathbf{d}_i-\mathbf{m}_i)^\top
    \mathbf{C}_i^{-1}(\mathbf{d}_i-\mathbf{m}_i)
  \right],
  \label{eq:loglike}
\end{equation}
with $k = N_{\rm bin}$ and $\mathbf{C}_i$ the per-cluster covariance.  The
implementation, \texttt{cholog\_gaussian}, Cholesky-factorises $\mathbf{C}_i$
once and evaluates $\ln\det$ and quadratic form via triangular solves.

\subsection{Forward model with fixed-parameter miscentering}

The forward model uses the parametric concentration of
Eq.~(\ref{eq:parametric_c}) and feeds $(M, c, z)$ to
\texttt{xlensing.model.NFW\_Delta\_Sigma}.  The returned dictionary contains
the four one-halo components $\{$\texttt{NFW Signal},
\texttt{Miscentered Signal}, \texttt{BCG Signal}, \texttt{Two-halo term},
\texttt{Signal}$\}$ with the mixing convention
\begin{equation}
  \texttt{Signal} = (1 - \mathrm{FMISS})\,\texttt{NFW Signal}
      + \mathrm{FMISS}\,\texttt{Miscentered Signal}
      + \texttt{BCG Signal}
      + \texttt{Two-halo term}.
  \label{eq:mixing}
\end{equation}
Our forward model returns
\begin{equation}
  \mathbf{m}_i = \texttt{Signal}_i - \texttt{Two-halo term}_i,
  \label{eq:model}
\end{equation}
i.e.\ the centred + miscentred NFW mixture, with the 2-halo term excluded
because the simulation contains only the 1-halo triaxial NFW signal.  The
miscentering parameters are held \emph{fixed} at the population values used to
generate the offsets, uniformly across the sample:
\begin{equation}
  \mathrm{FMISS}_i = 1 - \fcen = 0.36,\qquad
  \sigmaoff,\,i = 0.32\,h^{-1}\,\mathrm{Mpc} \quad\forall\,i.
\end{equation}
No per-cluster knowledge of which clusters are miscentred is exposed to the
fit—the model averages the Rayleigh+delta mixture identically over every
cluster.

\subsection{Triaxiality mitigation: radial cut}
\label{sec:radial_cut}

Spherical NFW fits to data simulated with triaxial NFW shears systematically
bias the recovered concentration normalisation $c_0$ low, and via the
$c_0$--$\gamma$ degeneracy bias $\gamma$ high.  The mismatch is largest in
the inner profile $R\lesssim r_s$ where the NFW peak is most shape-sensitive;
at $R \gtrsim r_s$ the $1/R^2$ outer tail is nearly shape-independent.

The notebook implements a model-independent diagnostic (described in
Section~\ref{sec:plots}) that motivates the analysis-side cut
\begin{equation}
  R_{\rm fit}^{(b)} \ge \Rmin = 0.5\,\mathrm{Mpc},
\end{equation}
applied by slicing the data vector, covariance, and radii array passed to
\texttt{run\_sampler\_combined}.  The simulation, forward model, and sampler
are otherwise unchanged.

\subsection{Priors}
\label{sec:priors}

Priors are declared explicitly as a dictionary of frozen
\texttt{scipy.stats} distributions and passed to
\texttt{run\_sampler\_combined} via the \texttt{priors} keyword:
\begin{equation}
  \begin{aligned}
    c_0          &\sim \mathcal{N}( 5.0, 0.5),  & \alpha       &\sim \mathcal{N}(-0.1, 0.05),\\
    \beta        &\sim \mathcal{N}(-0.4, 0.2),  & \gamma       &\sim \mathcal{N}( 0.0, 1.0),\\
    \log M_{\rm norm} &\sim \mathcal{N}( 0.0, 0.1). & &
  \end{aligned}
  \label{eq:priors}
\end{equation}
The priors on $c_0$, $\alpha$, $\beta$ are deliberately narrow Gaussians
centred on the Dutton--Maccio (2014) power-law expectation
(Eq.~\ref{eq:DM14}): the working assumption is that the $M$--$c$ relation is
\emph{pre-determined} from external calibration (simulations, complementary
data) and these parameters are sampled only to absorb small departures induced
by a non-zero $\gamma$.  The $\gamma$ prior is broad, $\mathcal{N}(0, 1)$, so
its posterior contraction provides a clean indication of magnitude-gap
information in the lensing data.

\subsection{Nested sampling}

Posteriors are obtained with \textsc{nautilus}
\citep{nautilus2023}, a flow-trained importance sampler, using
$n_{\rm live}=400$ live points and \texttt{multiprocess.Pool(15)}.  The
posterior is extracted as \texttt{sampler.posterior()} and reweighted to a
probability vector via $w \propto \exp(\log w - \max\log w)$.

%──────────────────────────────────────────────────────────────────────────────
\section{Notebook outputs}
\label{sec:plots}
%──────────────────────────────────────────────────────────────────────────────

In addition to the simulation-stage distributions already placed in
Section~2 (Figures \ref{fig:nz}, \ref{fig:truth}, \ref{fig:mc_scatter},
\ref{fig:triaxial}, \ref{fig:mis}), the notebook produces a further five
figures from the lensing measurement and the inference.  The three result
figures are saved to disk; the remaining two are inline diagnostics.

\subsection{Sample $\dsig$ profiles}
Six representative profiles, error bars and symbols colour-coded by $\dm$ and
shaped (\textcircled{}/\Large$\triangle$\normalsize) for centred/miscentred clusters.
The miscentred profiles visibly suppress the inner peak.

\subsection{Triaxiality residual diagnostic}
Two panels.  \emph{Left}: the fractional per-cluster residual
$(\dsig_{\rm data} - \dsig_{\rm NFW,truth})/\dsig_{\rm NFW,truth}$, where
$\dsig_{\rm NFW,truth}$ is the \emph{spherical-NFW} prediction evaluated at the
cluster's true $(M, c, z)$.  No fitting is involved—the scatter directly
measures the shape mismatch.  The interquartile band is broad at small $R$
and narrows above $\sim\!0.5$\,Mpc.  \emph{Right}: the half 16--84 percentile
fractional spread $\sigma_{\rm res}(R)$, showing that the dominant systematic
is concentrated at $R\lesssim 0.3$--$0.5$\,Mpc.  A vertical dashed line marks
$\Rmin = 0.5\,\mathrm{Mpc}$.

\subsection{Corner plot (saved)}
\begin{figure}[ht!]
  \centering
  \includegraphics[width=0.8\textwidth]{notebooks/combined_posterior_triaxial_miscentering_unbiased}
  \caption{Posterior corner plot with prior overlay.  Off-diagonal panels show
  2-D contours.  Diagonal panels (1-D marginals) overlay the prior PDF (dashed
  \textsf{darkorange}), the OLS truth value (\textsf{crimson}), and the
  density-normalised posterior histogram (\textsf{steelblue}).  Generated by
  \texttt{notebooks/combined\_recovery\_lensing\_triaxial\_unbiased.ipynb}.}
  \label{fig:corner}
\end{figure}
Figure~\ref{fig:corner} is the headline output.  With tight DM14 priors the
$c_0$, $\alpha$, $\beta$ diagonals essentially trace the priors (low data
contribution), as expected from our assumption of a pre-determined
$M$--$c$.  The $\gamma$ diagonal is strongly contracted relative to its
$\mathcal{N}(0, 1)$ prior, demonstrating that the magnitude-gap signal is
recovered from the lensing data.  The truth markers fall within the 68\% CIs.

\subsection{$M$--$c$ recovery band (saved)}
\begin{figure}[hb!]
  \centering
  \includegraphics[width=0.8\textwidth]{notebooks/combined_mc_recovery_triaxial_miscentering_unbiased}
  \caption{$M$--$c$ recovery band evaluated at $z=z_{\rm med}$ for three
  representative magnitude gaps $\dm \in \{0.5, \dm_{\rm ref}, 2.0\}$~mag.
  Filled bands show the 300-sample posterior envelope; solid lines are
  posterior medians; dashed curves are the DM14 truth curves at the same
  $\dm$.  Generated by the notebook.}
  \label{fig:mc}
\end{figure}
Figure~\ref{fig:mc} visualises the recovered $M$--$c$ relation at three
magnitude gaps.  The $\gamma\log_{10}(\dm/\dm_{\rm ref})$ term tilts the
relation upward (downward) for larger (smaller) $\dm$ relative to the
reference value.  The posterior medians track the DM14 truth curves modulated
by $\gamma$.

\subsection{$\gamma$ marginal posterior (saved)}
\begin{figure}[ht!]
  \centering
  \includegraphics[width=0.8\textwidth]{notebooks/combined_gamma_marginal_triaxial_miscentering_unbiased}
  \caption{Marginal posterior on the magnitude-gap slope $\gamma$.  The
  density-normalised histogram (\textsf{steelblue}) is the posterior, the
  shaded band is the 68\% credible interval, the solid vertical line is the
  median, the dashed \textsf{crimson} line is the truth, and the dashed
  \textsf{darkorange} curve is the $\mathcal{N}(0, 1)$ prior.  Generated by
  the notebook.}
  \label{fig:gamma}
\end{figure}
Figure~\ref{fig:gamma} provides the cleanest view of the $\gamma$ recovery:
the posterior is dramatically tighter than the prior, indicating substantial
information in the lensing data, and the median is within $\sim\!1\sigma$ of
the truth.

%──────────────────────────────────────────────────────────────────────────────
\section{Suggested actions before DES Y3 application}
\label{sec:future}
%──────────────────────────────────────────────────────────────────────────────

The pipeline is mature enough to constrain $\gamma$ on a self-consistent
DES~Y3-like mock with five-parameter inference.  Before it is applied to real
DES~Y3 data, a number of additions and validations should be carried out, in
roughly increasing order of scope.

\subsection{Source-side realism}

\subsubsection*{Tomographic source $n(z)$ with calibration uncertainty}

The Smail--Schechter form is a single-component approximation; the actual
DES~Y3 source distribution is the SOMPZ-calibrated, tomographic
$n(z)$ of \citet{Myles2021,MacCrann2021}.  Replace
\texttt{sample\_des\_y3\_source\_z} with a loader that reads the official
\texttt{des-y3-redshifts.hdf5} per-bin $n(z)$ tabulations and:
\begin{itemize}[noitemsep]
  \item Allows weighting per source by the relevant tomographic bin.
  \item Optionally marginalises over the SOMPZ shift and stretch parameters
        $(\Delta z_i, s_i)$ supplied as Gaussian nuisance priors.
\end{itemize}

\subsubsection*{Photometric redshift dilution}

Real photo-$z$s introduce a bias and dilution of the lensing signal that
must be propagated explicitly:
\begin{equation}
  \dsig_{\rm obs}(R) = \int p(z_s\,|\,z_{\rm phot})\,\dsig(R; z_l, z_s)\,\Sigma_{\rm cr}^{-1}(z_l,z_s)\,\mathrm{d}z_s,
\end{equation}
either by integrating per-source over $p(z_s|z_{\rm phot})$ or by computing the
bin-averaged $\langle\Sigma_{\rm cr}^{-1}\rangle$.

\subsubsection*{Shear calibration}

The mock uses an ideal lognormal-size metacalibration response; DES~Y3
provides per-source response matrices and selection responses.  These should
be propagated into the per-source weights and into the boost factor.

\subsubsection*{Boost factor / cluster-member contamination}

Cluster members contaminate the source sample and dilute the signal as a
function of radius:
\begin{equation}
  B(R) \equiv \frac{n_{\rm cl}(R)}{n_{\rm field}(R)} \to 1 \quad \mathrm{at\ large}\ R.
\end{equation}
A reasonable starting point is the parametric form of \citet{Varga2019} fit to
DES~Y3 stacks, applied either as a multiplicative correction
$\dsig_{\rm corr}(R) = B(R)\,\dsig_{\rm meas}(R)$ or by absorbing it as a free
nuisance.

\subsection{Cluster-side realism}

\subsubsection*{Real redMaPPer Y3 catalogue ingestion}

Replace the synthetic truth catalogue with a loader for the DES~Y3 GOLD
2.2.1 redMaPPer cluster catalogue, including:
\begin{itemize}[noitemsep]
  \item $(\mathrm{RA}, \mathrm{Dec}, z_\lambda, \lambda)$ columns.
  \item Per-cluster $P_{\rm CEN}$ (centring probability of the most-likely
        central candidate) so that the binary
        $\fcen$ can be replaced with a per-cluster $\fcen,i$
        \citep{Rykoff2016}.
  \item Membership probabilities $p_{\rm mem}$ used to construct the magnitude
        gap.
\end{itemize}

\subsubsection*{Magnitude gap from real photometry}

The HOD + power-law luminosity-function construction of $\dm$ should be
replaced by a $\dm$ measured from the red-sequence members of redMaPPer.
A pre-processing script should:
\begin{enumerate}[noitemsep]
  \item Identify candidate BCGs and second-brightest red members within
        $0.5\,R_\lambda$.
  \item Apply magnitude limits and membership-probability cuts.
  \item Propagate the photometric uncertainty to a per-cluster $\sigma_{\dm}$
        that can enter the likelihood as a Gaussian.
\end{enumerate}

\subsubsection*{Mass--richness relation with intrinsic scatter}

The current pipeline inverts the deterministic McClintock+18
\citep{McClintock2018} relation to map $\lambda \to M$.  Real-data use requires:
\begin{itemize}[noitemsep]
  \item Adding the lognormal intrinsic scatter $\sigma_{\ln M | \lambda}$ that
        the McClintock+18 likelihood marginalises over.
  \item Updating to the DES~Y3 mass--richness calibration if/when published
        for the Y3 catalogue.
\end{itemize}

\subsection{Forward-model improvements}

\subsubsection*{Triaxial-aware forward model}

The current radial cut at $\Rmin=0.5\,\mathrm{Mpc}$ is a pragmatic alternative
to modelling triaxiality.  A more rigorous approach is to use an
orientation-averaged triaxial NFW profile pre-computed for the assumed CDM-like
shape distribution \citep{Oguri2012, Corless2007, JingSuto2002}.  Once
implemented, the radial cut can be relaxed and the inner-profile information
on the concentration can be exploited.

\subsubsection*{2-halo term}

When fitting real data we cannot subtract the 2-halo term as we do here.
Either:
\begin{itemize}[noitemsep]
  \item Include the 2-halo term in the forward model
        (\texttt{xlensing.model.NFW\_Delta\_Sigma['Two-halo term']} or an
        explicit Eisenstein--Hu power-spectrum-based prediction), or
  \item Restrict the maximum radius of the fit to $R_{\rm max} \lesssim
        4$--$5\,\mathrm{Mpc}$, where the 2-halo contribution remains small.
\end{itemize}

\subsubsection*{Per-cluster miscentering from $P_{\rm CEN}$}

Section~\ref{sec:miscentering_sim} fixes
$\mathrm{FMISS}=1-\fcen$ uniformly.  For real data the redMaPPer
$P_{\rm CEN}$ probability per cluster is available; this can be promoted from
a population constant to a per-cluster prior:
\begin{equation}
  \mathrm{FMISS}_i \sim \mathrm{Beta}(\alpha_{P_{\rm CEN},i}, \beta_{P_{\rm CEN},i})
  \quad\mathrm{or}\quad \mathrm{FMISS}_i = 1 - P_{{\rm CEN},i}.
\end{equation}

\subsection{Validation tests}

\subsubsection*{Recovery of a $\gamma = 0$ null}

Run the same pipeline with $\gamma_{\rm true} = 0$ and verify that the
posterior on $\gamma$ contracts to $0$ within errors.  This is the headline
null test.

\subsubsection*{Recovery vs $N$}

Vary $N \in \{1, 5, 10, 16, 25\}\times 10^3$ and verify
$\sigma_\gamma \propto 1/\sqrt{N}$ and that biases scale at most logarithmically
with $N$.

\subsubsection*{Recovery vs $\Rmin$}

Sweep $\Rmin$ over $\{0.2, 0.3, 0.5, 0.75, 1.0\}\,\mathrm{Mpc}$ and verify that
the bias on $(c_0, \gamma)$ plateaus near zero for $\Rmin \gtrsim r_s$ and that
the posterior widths grow predictably.

\subsubsection*{Sensitivity to miscentering misspecification}

Repeat the analysis after deliberately misspecifying the model:
\begin{itemize}[noitemsep]
  \item Set $\mathrm{FMISS}_{\rm model} = 0$ while the simulation has
        $\fcen=0.64$.  Quantify the resulting bias on $(c_0, \gamma)$.
  \item Set $\sigmaoff^{\rm model} = (0.16, 0.48)\,h^{-1}\mathrm{Mpc}$ while
        the simulation uses $0.32\,h^{-1}\mathrm{Mpc}$.  This bounds the
        required prior width on $\sigmaoff$ in a future free-parameter
        version of the fit.
\end{itemize}

\subsubsection*{Stacked vs destacked comparison}

Bin the same mock data in $\lambda$ as in McClintock+18, run a stacked NFW
fit per bin, and compare the resulting $\sigma_{c_0}$ to the destacked
result.  The prototype showed a factor $\sim\!2$ improvement; this should be
reproduced under the more realistic DES Y3-like setup.

\subsection{Documentation and reproducibility}

\begin{itemize}[noitemsep]
  \item Seed every RNG so the validation tests above are exactly reproducible.
  \item Persist the truth catalogue, lensing data, covariance, and sampler
        checkpoint to disk so re-analysis with a different forward model
        does not require re-running the lensing simulation.
  \item Pin \texttt{xlensing}, \texttt{nautilus}, \texttt{colossus},
        \texttt{ngmix} versions in \texttt{pyproject.toml}.
  \item Add a regression test that runs a small ($N=200$) version of the
        notebook on every push.
\end{itemize}

%──────────────────────────────────────────────────────────────────────────────
\section{Conclusions}
%──────────────────────────────────────────────────────────────────────────────

The notebook
\texttt{combined\_recovery\_lensing\_triaxial\_unbiased.ipynb} demonstrates
the joint un-stacked inference of the $M$--$c$ relation parameters
$(c_0, \alpha, \beta)$ and the magnitude-gap slope $\gamma$ from per-cluster
$\dsig$ profiles in a setup that contains the three dominant systematics
relevant to a DES~Y3 cluster lensing analysis:
\begin{enumerate}[noitemsep]
  \item triaxial halo geometry,
  \item Rayleigh + delta-function miscentering, and
  \item realistic redshift distributions for both lenses and sources.
\end{enumerate}

The recovery is unbiased once the analysis is configured with
(i)~tight Gaussian priors on the $M$--$c$ parameters centred on the DM14
expectation, (ii)~a broad $\mathcal{N}(0, 1)$ prior on the magnitude-gap slope
$\gamma$, (iii)~a fixed-parameter Rayleigh miscentering model in the forward
pass with $\mathrm{FMISS}=1-\fcen$ and the same $\sigmaoff$ used in the
simulation, and (iv)~an analysis-side radial cut at $\Rmin = 0.5\,\mathrm{Mpc}$
that suppresses the inner-bin triaxiality mismatch.

The list of Section~\ref{sec:future} enumerates the work required before the
pipeline is applied to real data.  The largest remaining items are the
ingestion of the actual DES~Y3 source $n(z)$ and shear catalogue, the
inclusion of photometric-redshift dilution and boost-factor correction,
and—on the modelling side—a triaxial-aware forward model that would allow
the radial cut to be relaxed.

%──────────────────────────────────────────────────────────────────────────────
\begin{thebibliography}{99}

\bibitem[Allgood et al.(2006)]{Allgood2006}
Allgood, B., Flores, R.~A., Primack, J.~R., et al.\ 2006, \mnras, 367, 1781

\bibitem[Corless \& King(2007)]{Corless2007}
Corless, V.~L., \& King, L.~J.\ 2007, \mnras, 380, 149

\bibitem[Costanzi et al.(2021)]{Costanzi2021}
Costanzi, M., Saro, A., Bocquet, S., et al.\ 2021, \prd, 103, 043522

\bibitem[Diemer \& Joyce(2019)]{Diemer2019}
Diemer, B., \& Joyce, M.\ 2019, \apj, 871, 168

\bibitem[Dutton \& Macci\`o(2014)]{DuttonMaccio2014}
Dutton, A.~A., \& Macci\`o, A.~V.\ 2014, \mnras, 441, 3359

\bibitem[Jing \& Suto(2002)]{JingSuto2002}
Jing, Y.~P., \& Suto, Y.\ 2002, \apj, 574, 538

\bibitem[Jones et al.(2003)]{Jones2003}
Jones, L.~R., Ponman, T.~J., Horton, A., et al.\ 2003, \mnras, 343, 627

\bibitem[Kawahara(2010)]{Kawahara2010}
Kawahara, H.\ 2010, \apj, 719, 1926

\bibitem[Khosroshahi et al.(2006)]{Khosroshahi2006}
Khosroshahi, H.~G., Jones, L.~R., \& Ponman, T.~J.\ 2006, \mnras, 372, L68

\bibitem[Lange(2023)]{nautilus2023}
Lange, J.~U.\ 2023, \mnras, 525, 3181

\bibitem[MacCrann et al.(2021)]{MacCrann2021}
MacCrann, N., Becker, M.~R., McCullough, J., et al.\ 2021, \mnras, 509, 3371

\bibitem[McClintock et al.(2018)]{McClintock2018}
McClintock, T., Varga, T.~N., Gruen, D., et al.\ 2018, \mnras, 482, 1352

\bibitem[Myles et al.(2021)]{Myles2021}
Myles, J., Alarcon, A., Amon, A., et al.\ 2021, \mnras, 505, 4249

\bibitem[Oguri \& Hamana(2011)]{Oguri2012}
Oguri, M., \& Hamana, T.\ 2011, \mnras, 414, 1851

\bibitem[Rykoff et al.(2016)]{Rykoff2016}
Rykoff, E.~S., Rozo, E., Hollowood, D., et al.\ 2016, \apjs, 224, 1

\bibitem[Schneider et al.(2012)]{Schneider2012}
Schneider, M.~D., Frenk, C.~S., \& Cole, S.\ 2012, \jcap, 2012, 030

\bibitem[Sereno et al.(2018)]{Sereno2018}
Sereno, M., Umetsu, K., Ettori, S., et al.\ 2018, \mnrasl, 473, L89

\bibitem[Smail et al.(1994)]{Smail1994}
Smail, I., Ellis, R.~S., \& Fitchett, M.~J.\ 1994, \mnras, 270, 245

\bibitem[Tinker et al.(2008)]{Tinker2008}
Tinker, J., Kravtsov, A.~V., Klypin, A., et al.\ 2008, \apj, 688, 709

\bibitem[Varga et al.(2019)]{Varga2019}
Varga, T.~N., DeRose, J., Gruen, D., et al.\ 2019, \mnras, 489, 2511

\bibitem[Vitorelli et al.(2025)]{Vitorelli2025}
Vitorelli, A., \textit{destacking: joint un-stacked fitting of cluster
weak-lensing profiles}, \url{https://github.com/andrevitorelli/destacking}, 2025.

\bibitem[Wechsler et al.(2002)]{Wechsler2002}
Wechsler, R.~H., Bullock, J.~S., Primack, J.~R., Kravtsov, A.~V., \& Dekel, A.\ 2002, \apj, 568, 52

\bibitem[Vitorelli(xlensing)]{xlensing}
Vitorelli, A., \textit{xlensing: weak-lensing analysis tools for galaxy clusters},
\url{https://github.com/andrevitorelli/xlensing}.

\bibitem[Zhang et al.(2019)]{Zhang2019}
Zhang, Y., Jeltema, T., Hollowood, D.~L., et al.\ 2019, \mnras, 487, 2578

\bibitem[Zheng et al.(2007)]{Zheng2007}
Zheng, Z., Coil, A.~L., \& Zehavi, I.\ 2007, \apj, 667, 760

\end{thebibliography}

\end{document}
